\documentclass[12pt, a4paper]{article}

% ── Codificación y fuentes ──────────────────────────────────────────────────
\usepackage[utf8]{inputenc}
\usepackage[T1]{fontenc}
\usepackage{lmodern}

% ── Matemáticas ─────────────────────────────────────────────────────────────
\usepackage{amsmath}
\usepackage{amssymb}
\usepackage{mathtools}          % coloneqq, aligned, etc.

% ── Idioma ──────────────────────────────────────────────────────────────────
\usepackage[spanish]{babel}

% ── Diseño de página ────────────────────────────────────────────────────────
\usepackage[top=2.5cm, bottom=2.5cm, left=3cm, right=3cm]{geometry}
\usepackage{parskip}            % espacio entre párrafos en lugar de sangría
\usepackage{microtype}          % mejor justificación del texto

% ── Colores y cajas ─────────────────────────────────────────────────────────
\usepackage{xcolor}
\usepackage{tcolorbox}
\tcbuselibrary{skins, breakable}

% ── Encabezados y secciones ─────────────────────────────────────────────────
\usepackage{titlesec}
\titleformat{\section}{\large\bfseries\color{blue!60!black}}{}{0em}{}[\titlerule]
\titlespacing{\section}{0pt}{1.5ex}{0.8ex}

% ── Enumeraciones y listas ──────────────────────────────────────────────────
\usepackage{enumitem}
\setlist[itemize]{noitemsep, topsep=3pt}

% ── Hiperreferencias ────────────────────────────────────────────────────────
\usepackage{hyperref}
\hypersetup{colorlinks=true, linkcolor=blue!60!black}

% ── Macros de conveniencia ──────────────────────────────────────────────────
\newcommand{\Lap}{\mathcal{L}}              % Transformada de Laplace
\newcommand{\ILap}{\mathcal{L}^{-1}}       % Transformada inversa
\newcommand{\Uu}{\mathcal{U}}              % Función escalón de Heaviside
\newcommand{\half}{\tfrac{1}{2}}
\newcommand{\sfrac}[2]{\tfrac{#1}{#2}}

% ── Definición de caja para resultados destacados ───────────────────────────
\newtcolorbox{resultbox}{
  enhanced, breakable,
  colback=blue!5!white, colframe=blue!60!black,
  fonttitle=\bfseries, title=Resultado,
  left=6pt, right=6pt, top=4pt, bottom=4pt
}

\newtcolorbox{stepbox}[1]{
  enhanced,
  colback=gray!8!white, colframe=gray!50!black,
  title=#1, fonttitle=\bfseries\small,
  left=6pt, right=6pt, top=4pt, bottom=4pt
}

% ── Datos del documento ─────────────────────────────────────────────────────
\title{\textbf{Solución mediante Transformada de Laplace}\\[4pt]
       \large Ecuación diferencial de segundo orden con forzamiento discontinuo}
\date{}

% ════════════════════════════════════════════════════════════════════════════
\begin{document}
\maketitle
\thispagestyle{empty}

% ────────────────────────────────────────────────────────────────────────────
\section{Planteamiento del problema}

Se resuelve el problema de valor inicial:
\begin{equation}
  y'' + y' + \frac{5}{4}\,y = 1 - \Uu(t-\pi),
  \qquad y(0) = 1,\quad y'(0) = -1.
  \label{eq:pvi}
\end{equation}

% ────────────────────────────────────────────────────────────────────────────
\section{Aplicación de la Transformada de Laplace}

Sea $Y = \Lap\{y\}$. Las transformadas de las derivadas son:
\begin{align*}
  \Lap\{y'\}  &= sY - y(0)         = sY - 1, \\
  \Lap\{y''\} &= s^{2}Y - sy(0) - y'(0) = s^{2}Y - s + 1.
\end{align*}

Para el lado derecho, usando el segundo teorema de traslación:
\[
  \Lap\!\left\{1 - \Uu(t-\pi)\right\} = \frac{1}{s} - \frac{e^{-\pi s}}{s}.
\]

Sustituyendo en \eqref{eq:pvi}:
\[
  \bigl(s^{2}Y - s + 1\bigr) + \bigl(sY - 1\bigr) + \frac{5}{4}\,Y
  = \frac{1}{s} - \frac{e^{-\pi s}}{s},
\]
\[
  Y\!\left(s^{2} + s + \frac{5}{4}\right) = \frac{1}{s} - \frac{e^{-\pi s}}{s} + s.
\]

\paragraph{Factorización del denominador.}
Completando el cuadrado:
\[
  s^{2} + s + \frac{5}{4}
  = \left(s + \half\right)^{2} + 1.
\]

Por lo tanto:
\begin{equation}
  Y = \underbrace{\frac{1}{s\!\left[\!\left(s+\half\right)^{2}+1\right]}}_{Y_1}
    - \underbrace{\frac{e^{-\pi s}}{s\!\left[\!\left(s+\half\right)^{2}+1\right]}}_{Y_2}
    + \underbrace{\frac{s}{\left(s+\half\right)^{2}+1}}_{Y_3}.
  \label{eq:Y}
\end{equation}

% ────────────────────────────────────────────────────────────────────────────
\section{Cálculo de las transformadas inversas}

% ---- L1 --------------------------------------------------------------------
\begin{stepbox}{Transformada $L_1 = \ILap\{Y_1\}$: convolución}

Escribimos $Y_1 = F(s)\cdot G(s)$ con
\[
  F(s) = \frac{1}{s} \;\Longrightarrow\; f(t) = 1,
  \qquad
  G(s) = \frac{1}{\left(s+\half\right)^{2}+1}
         \;\Longrightarrow\; g(t) = e^{-t/2}\sin t.
\]

Por el teorema de convolución:
\[
  L_1 = (f * g)(t) = \int_{0}^{t} e^{-\tau/2}\sin\tau\,d\tau.
\]

\textbf{Integral indefinida} (integración por partes dos veces):

Primera aplicación con $u = e^{-\tau/2}$, $dv = \sin\tau\,d\tau$:
\[
  \int e^{-\tau/2}\sin\tau\,d\tau
  = -\cos\tau\,e^{-\tau/2}
    - \frac{1}{2}\int \cos\tau\,e^{-\tau/2}\,d\tau.
\]

Segunda aplicación con $u = e^{-\tau/2}$, $dv = \cos\tau\,d\tau$:
\[
  \int e^{-\tau/2}\sin\tau\,d\tau
  = -\cos\tau\,e^{-\tau/2}
    - \frac{1}{2}\!\left(\sin\tau\,e^{-\tau/2}
    - \frac{1}{2}\int e^{-\tau/2}\sin\tau\,d\tau\right).
\]

Despejando la integral:
\[
  \frac{5}{4}\int e^{-\tau/2}\sin\tau\,d\tau
  = -\cos\tau\,e^{-\tau/2} - \frac{1}{2}\sin\tau\,e^{-\tau/2},
\]
\[
  \int e^{-\tau/2}\sin\tau\,d\tau
  = -\frac{4}{5}\cos\tau\,e^{-\tau/2} - \frac{2}{5}\sin\tau\,e^{-\tau/2}.
\]

Evaluando en $[0,\,t]$:
\end{stepbox}

\begin{resultbox}
\[
  L_1 = \frac{4}{5} - \frac{4}{5}\cos t\,e^{-t/2} - \frac{2}{5}\sin t\,e^{-t/2}.
\]
\end{resultbox}

\bigskip

% ---- L2 --------------------------------------------------------------------
\begin{stepbox}{Transformada $L_2 = \ILap\{Y_2\}$: segundo teorema de traslación}

Por el segundo teorema de traslación, si $\ILap\!\left\{\sfrac{F(s)}{s}\right\} = f(t)$, entonces:
\[
  \ILap\!\left\{e^{-\pi s}\frac{F(s)}{s}\right\} = \Uu(t-\pi)\,f(t-\pi).
\]

Usando el resultado de $L_1$ con $t \leftarrow t - \pi$:
\end{stepbox}

\begin{resultbox}
\[
  L_2 = \Uu(t-\pi)\!\left[
      \frac{4}{5}
    - \frac{4}{5}\cos(t-\pi)\,e^{-(t-\pi)/2}
    - \frac{2}{5}\sin(t-\pi)\,e^{-(t-\pi)/2}
  \right].
\]
\end{resultbox}

\bigskip

% ---- L3 --------------------------------------------------------------------
\begin{stepbox}{Transformada $L_3 = \ILap\{Y_3\}$: traslación en $s$}

Sumando y restando $\half$ en el numerador:
\[
  Y_3 = \frac{s + \half - \half}{\left(s+\half\right)^{2}+1}
      = \frac{s + \half}{\left(s+\half\right)^{2}+1}
        - \frac{\half}{\left(s+\half\right)^{2}+1}.
\]

Aplicando la primera propiedad de traslación ($e^{-at}\cos\omega t$ y $e^{-at}\sin\omega t$):
\end{stepbox}

\begin{resultbox}
\[
  L_3 = e^{-t/2}\cos t - \frac{1}{2}\,e^{-t/2}\sin t.
\]
\end{resultbox}

% ────────────────────────────────────────────────────────────────────────────
\section{Solución final}

La solución del PVI \eqref{eq:pvi} es $y = L_1 - L_2 + L_3$:

\begin{tcolorbox}[enhanced, colback=green!6!white, colframe=green!50!black,
                  fonttitle=\bfseries, title=Solución $y(t)$]
\begin{align*}
  y(t) ={}&
    \frac{4}{5}
    - \frac{4}{5}\cos t\,e^{-t/2}
    - \frac{2}{5}\sin t\,e^{-t/2}
    + e^{-t/2}\cos t
    - \frac{1}{2}\,e^{-t/2}\sin t \\[6pt]
    &- \Uu(t-\pi)\!\left[
        \frac{4}{5}
        - \frac{4}{5}\cos(t-\pi)\,e^{-(t-\pi)/2}
        - \frac{2}{5}\sin(t-\pi)\,e^{-(t-\pi)/2}
      \right].
\end{align*}
\end{tcolorbox}

\end{document}